\section{Rice's theorem}
\label{sec:universal-reductions-05}

\begin{definition}[Extensional index property]
A set \(A\subseteq\N\) is \emph{extensional} iff
\[
  \varphi_e\simeq\varphi_d
  \quad\Longrightarrow\quad
  (e\in A\Longleftrightarrow d\in A)
  \qquad\text{for all }e,d\in\N.
\]
It is \emph{nontrivial} iff it is neither empty nor all of \(\N\).
\end{definition}

\begin{theorem}[Rice]
\label{thm:rice}
Every nontrivial extensional set of indices of unary partial computable
functions is undecidable.
\textit{Status: complete proof.}
\end{theorem}

\begin{proof}
Let \(A\) be nontrivial and extensional.  First suppose an index of the
everywhere-undefined function does not belong to \(A\).  Choose \(a\in A\).
Fix an index \(b\) for the binary partial computable function
\[
 F(e,x)\simeq
 \begin{cases}
   \varphi_a(x),&\text{if }\varphi_e(e)\text{ returns a value},\\
   \text{undefined},&\text{otherwise}.
 \end{cases}
\]
It is partial computable by first simulating \(U(e,e)\) and, only after that
simulation returns, simulating \(U(a,x)\); \cref{thm:enumeration} then
supplies the index \(b\).  Define the total computable map
\(g(e)=\mathsf{par}_1^1(b,e)\).  The parameter theorem gives
\(\varphi_{g(e)}(x)\simeq F(e,x)\) for every \(x\).  If \(e\notin K\), the
new program is everywhere undefined, so \(g(e)\notin A\).  If \(e\in K\), the
new program computes \(\varphi_a\), so \(g(e)\in A\).  Hence
\(K\leq_m A\), and \(A\) is undecidable.

If the everywhere-undefined function has an index in \(A\), apply the first
case to the nontrivial extensional complement \(\N\setminus A\).  Decidability
of \(A\) would imply decidability of its complement, which is impossible.
\end{proof}

\begin{warning}
Rice's theorem concerns semantic properties of the partial function computed
by a program.  Syntactic properties such as ``the source contains ten states''
can be decidable and are not extensional.
\end{warning}

